\metatitle{Recurrence limit sets and Toeplitz minors}
\metadescription{Limit sets of zeros of polynomial recurrences, the Beraha--Kahane--Weiss theorem, renewal families, Toeplitz minors, the Schmidt--Spitzer theorem, and real spectra.}
\metakeywords{recurrence limit set, limiting zeros, Toeplitz determinant, Toeplitz minor, Beraha--Kahane--Weiss, Schmidt--Spitzer, limiting spectrum}

\section[toeplitzDeterminants]{Recurrence limit sets and Toeplitz minors}

The zeros of a \hyperref[polynomialRecurrences]{polynomial recurrence}
often accumulate on curves or intervals. We study these limit sets through
characteristic roots, then specialize to Toeplitz determinants and minors.
The \hyperref[toeplitzBkwTheorem]{Beraha--Kahane--Weiss theorem}
gives the general analytic criterion; the
\hyperref[toeplitzSchmidtSpitzer]{Schmidt--Spitzer theorem}
describes the limiting spectra of banded Toeplitz sections.
The later finite-size proofs address the separate question of when
\emph{every} polynomial has only real zeros.

\subsection[recurrenceZeroLimitSets]{Limit sets and characteristic roots}

For a sequence of nonzero polynomials $P_n(x)$, its \defin{zero limit set} is
\[
 \mathcal L(P)=\{x\in\setC:
   x_j\to x,\ P_{n_j}(x_j)=0,\ n_j\to\infty
   \text{ for some sequences }(x_j),(n_j)\}.
\]
A finite-order recurrence whose coefficients are independent of $n$
has a generating function rational in the index variable $u$.
After canceling common factors, suppose its poles are simple on a parameter
domain. Partial fractions then give, for all sufficiently large $n$,
\[
 P_n(x)=\sum_i a_i(x)\lambda_i(x)^n,
\]
where $\lambda_i(x)$ are reciprocals of the poles.
The amplitudes $a_i$ depend on the initial conditions.
One must therefore consider both ties between dominant moduli
$|\lambda_i|$ and zeros of the corresponding amplitudes.
For the algebraic construction, including repeated poles, see the
\hyperref[recurrenceSpectralModel]{spectral expansion of a recurrence}.

\section[toeplitzBkw]{Beraha--Kahane--Weiss and limiting zeros}

The characteristic roots of a linear recurrence also control where zeros accumulate as
$n\to\infty$. This is the content of the Beraha--Kahane--Weiss theorem
\cite{Beraha75limits}, also used in the
\hyperref[berahaKahaneWeiss]{transfer-matrix setting}.
We now let $x$ be complex and work on a domain where the characteristic roots
and amplitudes are single-valued analytic functions.

\begin{theorem}[Beraha--Kahane--Weiss, analytic form]
\label{toeplitzBkwTheorem}
Let $d\geq2$ and let $a_i,\lambda_i$ be analytic on a connected open set
$\Omega\subseteq\setC$, with none identically zero. Suppose
$\lambda_i\not\equiv\omega\lambda_j$ for $i\ne j$ and every constant
$|\omega|=1$. Set
\[
 P_n(x)=\sum_{i=1}^d a_i(x)\lambda_i(x)^n.
\]
Call $i$ \defin{dominant} at $x$ if
$|\lambda_i(x)|=\max_j|\lambda_j(x)|$.
A point $x\in\Omega$ is a limit of zeros $x_k$ of $P_{n_k}$ with
$n_k\to\infty$ if and only if either
\begin{itemize}
\item exactly one index $i$ is dominant and $a_i(x)=0$; or
\item at least two indices are dominant.
\end{itemize}
Every neighborhood of such a point contains a zero of $P_n$ for all
sufficiently large $n$.
\end{theorem}

In \cite[Thm.~1.5]{Sokal2004ChromaticRoots}, the condition on constant
unit-modulus ratios is imposed on branches that are dominant on a set
with nonempty interior.

\begin{example}
For $P_n(x)=1+x^n+(x+1)^n$, all amplitudes equal $1$, so there are no
isolated limits arising from amplitude zeros. The limiting set is
\begin{align*}
 &\{x:|x|=1,\ |x+1|\leq1\}\\
 {}\cup{}&\{x:|x+1|=1,\ |x|\leq1\}\\
 {}\cup{}&\{x:|x|=|x+1|\geq1\}.
\end{align*}
The first two sets are circular arcs. The third consists of the two rays
\[
 \operatorname{Re}x=-\tfrac12,\qquad
 |\operatorname{Im}x|\geq\tfrac{\sqrt3}{2}.
\]
All three branches meet at $(-1\pm i\sqrt3)/2$. These equations follow by
comparing the moduli of $1,x,x+1$. In particular, the nonreal limiting points
force nonreal zeros for all sufficiently large $n$.
\end{example}

The \hyperref[firstDerivativeCriterion]{derivative-recurrence criterion} uses the values and logarithmic
derivatives of the characteristic roots; BKW uses their relative moduli.
Both depend on the amplitudes, which are determined by the initial polynomials.

For repeated poles the amplitudes $a_i(n,x)$ depend polynomially on $n$,
so the displayed theorem does not apply directly. A common scalar factor
in $n$ can be removed for large $n$; general polynomial-in-$n$ amplitudes
require an extension.

\subsection[recurrenceRenewalLimitSet]{Limiting zeros of the renewal family}

For the \hyperref[weightedRenewalInterlacing]{weighted renewal family},
BKW gives an explicit limiting interval.
The interlacing theorem applies to every index; the following statement
describes the zeros as the index tends to infinity.

\begin{theorem}[Renewal zero accumulation]
Let $d,\ell\geq1$ be integers and let
$R(u)=\prod_{j=1}^d(1+a_ju)$, with $a_j>0$.
Define $P_n$ by
\[
 \sum_{n\geq0}P_n(x)t^n=\frac1{1-tR(xt^\ell)}.
\]
There is a unique $u_*$ between the largest zero of $R$ and $0$ such that
\[
 R(u_*)+\ell u_*R'(u_*)=0.
\]
Put $x_*=u_*R(u_*)^\ell$. The set of finite accumulation points of zeros
of $P_n$, as $n\to\infty$, is exactly
\[
 (-\infty,x_*].
\]
Every neighborhood of every point in this ray contains a zero of $P_n$
for all sufficiently large $n$.
\end{theorem}

\begin{proof*}
Put $H(u)=uR(u)^\ell$ and $b=1/\max_j a_j$.
For $x\ne0$, poles of the generating function correspond bijectively to
\[
 H(u)=x,\qquad t=\frac1{R(u)},\qquad \lambda=R(u).
\]
Since $|\lambda|=(|x|/|u|)^{1/\ell}$, dominant characteristic roots
correspond to smallest-modulus roots of $H(u)=x$.
At a simple pole the amplitude is
\[
 A(u)=\frac{R(u)}{R(u)+\ell uR'(u)},
\]
which is nonzero.

For fixed $r>0$,
\[
 |H(re^{i\theta})|
 =r\prod_j(1+2a_jr\cos\theta+a_j^2r^2)^{\ell/2}
\]
is strictly decreasing for $0\leq\theta\leq\pi$.
Thus distinct roots of $H(u)=x$ with the same modulus must be conjugate,
and this can happen only for real $x$.
In particular, there is no identically unit-modulus ratio of distinct
characteristic branches, as required by BKW.

On $(-b,0)$, the function
\[
 g(u)=1+\ell\frac{uR'(u)}{R(u)}
 =1+\ell\sum_j\frac{a_ju}{1+a_ju}
\]
is strictly increasing from $-\infty$ to $1$.
Its unique zero is $u_*$, and $H''(u_*)=R(u_*)^\ell g'(u_*)>0$.
The inverse root near zero is uniquely smallest in modulus and remains so
for $x_*<x<0$. At $x_*$ it merges with another root; just below $x_*$
the two smallest roots form a conjugate pair.

This pair persists for all $x<x_*$.
Indeed, all zeros of $H'$ are real, and the upper-half-plane member satisfies
\[
 \arg u+\ell\sum_j\arg(1+a_ju)=\pi,
\]
with each argument in $(0,\pi)$.
It cannot encounter a critical point off the real axis.
A negative real limit beyond $-b$ would make the left side exceed $\pi$;
a limit in $(-b,0)$ would require a collision at $u_*$.
A positive real limit has positive $H$, and roots remain bounded over
bounded sets of $x$. Hence the pair continues along the whole ray.
It cannot lose dominance, since that would require a third root of the
same modulus.

Away from the finitely many critical values, BKW now gives accumulation
on the ray. Neighborhoods of the exceptional points contain regular
points of the ray, giving the same conclusion there, including at $x_*$.
On $(x_*,0)$ the unique dominant pole is simple with nonzero amplitude,
so its uniform coefficient asymptotics exclude accumulation, even if
subdominant poles collide. At $x=0$ the denominator is $1-t$, with a
single nearby dominant pole, so the same exclusion holds.
The renewal interlacing theorem excludes nonreal and positive zeros.
\end{proof*}

For the seed $R(u)=(1+u)^d$, put $L=d\ell+1$.
Then
\[
 u_*=-\frac1L,\qquad
 x_*=-\frac{(L-1)^{L-1}}{L^L}.
\]
For example, $L=2$ gives $(-\infty,-1/4]$, and $L=3$ gives
$(-\infty,-4/27]$.
Writing $P_n=C_{dn}$, where $\sum_{N\geq0}C_N(x)t^N=1/(1-t-xt^L)$,
identifies the same endpoint with the sign-reversed interval in
\cite[Thm.~1]{ForgcsTran2016} for the coefficient sequence $C_N$.
The dominance argument above also handles arbitrary positive row weights.

The \hyperref[weightedDelayRecurrences]{arbitrary-weight variable-predecessor theorem} need not come from a
stationary finite-order recurrence. BKW applies to the renewal specialization
above; more general changing coefficients require additional asymptotic
hypotheses.

\section[toeplitzMinors]{Toeplitz determinants and minors}

For a Laurent polynomial $b(z)=\sum_{k=-r}^s b_kz^k$, where
$r,s\geq1$ and $b_{-r}b_s\ne0$, the \defin{Toeplitz sections} are
\[
 T_n(b)=(b_{i-j})_{i,j=1}^n,\qquad
 P_n(t)=\det(tI-T_n(b)),\qquad P_0=1.
\]
The Laurent polynomial $b$ is the \defin{symbol}.
A \defin{Toeplitz minor} is the determinant of a square submatrix obtained
by selecting equally many rows and columns.

\subsection[toeplitzSchurMinors]{Schur functions as minors}

Let $h_m=\completeH_m(\xvec)$ be the complete homogeneous symmetric function,
with $h_0=1$ and $h_m=0$ for $m<0$.
For increasing nonnegative row indices $i_1<\dotsb<i_k$ and column indices
$j_1<\dotsb<j_k$ with $i_\ell\geq j_\ell$, put
\[
 \lambda_p=i_{k+1-p}+p-k,\qquad
 \mu_p=j_{k+1-p}+p-k\qquad(1\leq p\leq k).
\]
These are partitions with $\mu\subseteq\lambda$, and the
\hyperref[skewJacobiTrudi]{skew Jacobi--Trudi identity} gives
\[
 \det(h_{i_p-j_q})_{p,q=1}^k=\schurS_{\lambda/\mu}.
\]
Indeed, reversing both index orders gives the matrix
$(h_{\lambda_p-\mu_q-p+q})$ and leaves its determinant unchanged.
If some $i_\ell<j_\ell$, the lower-triangular zero pattern makes the minor zero.
The dual identity gives analogous formulas with elementary symmetric functions.

Nonnegative specializations therefore give
\hyperref[pfSchurJacobiTrudi]{totally non-negative Toeplitz matrices}.
For the general connection with coefficient sequences and real zeros, see
\hyperref[polyaFrequency]{Pólya frequency sequences};
for banded minors, determinant formulas, and recurrences, see
\cite{Alexandersson2012Toeplitz}.
The \hyperref[toeplitzReciprocalRecurrence]{rectangular Schur formula below}
makes the recurrence explicit for a balanced pentadiagonal matrix.


The \hyperref[toeplitzSchmidtSpitzer]{Schmidt--Spitzer theorem} below
describes limiting zeros of the principal characteristic determinants
$\det(tI-T_n(b))$. For other families of Toeplitz minors, one must identify
their own recurrence roots and amplitudes before applying BKW;
see \cite{Alexandersson2012Toeplitz}.

\section[toeplitzSpectra]{Limiting spectra of banded Toeplitz matrices}

\subsection[toeplitzSchmidtSpitzer]{The Schmidt--Spitzer theorem}

The \defin{limiting spectrum} is
\[
 \Lambda(b)=\{t\in\setC:\lim_{n\to\infty}
 \operatorname{dist}(t,\operatorname{spec}T_n(b))=0\}.
\]
\begin{theorem}[Schmidt--Spitzer, \cite{SchmidtSpitzer1960}]
For the banded symbol $b(z)=\sum_{k=-r}^s b_kz^k$ with
$r,s\geq1$ and $b_{-r}b_s\ne0$, order the $r+s$ roots of
$z^r(b(z)-t)$ by modulus, counting multiplicity:
\[
 |z_1(t)|\leq\dotsb\leq|z_{r+s}(t)|.
\]
Then
\[
 \Lambda(b)=\{t:|z_r(t)|=|z_{r+1}(t)|\}.
\]
\end{theorem}

Thus the limiting zeros of the principal characteristic determinants
$P_n(t)$ are given by an equal-modulus condition at the band index $r$.
See the \hyperref[toeplitzReciprocalRecurrence]{recurrence and BKW interpretation}
for the balanced pentadiagonal case.

\subsection[toeplitzRealLimitQuestion]{Real limiting spectra and finite spectra}

For real coefficients, \name{Boris Shapiro} and \name{František Štampach}
\cite{ShapiroStampach2019,ShapiroStampach2023Correction} considered three conditions:
all finite spectra are real; $\Lambda(b)\subseteq\setR$; and
$b^{-1}(\setR)$ contains a Jordan curve surrounding $0$.
The first implies the second, and the second implies the third.
The claimed converse involving the Jordan curve was withdrawn in the correction.
\name{Fabian Rohner} and \name{Clemens Thalhammer}
\cite[Thm.~1.3]{RohnerThalhammer2026} disprove even the implication from
the Jordan-curve condition to a real limiting spectrum, using the real symbol
\[
 b(z)=\left(z^2+97z+103-\frac{10201}{z}\right)^2.
\]
Their example does not settle whether a real limiting spectrum forces every
finite spectrum to be real. The following proposition settles this implication
for balanced pentadiagonal symbols with positive product of the outer coefficients.
The argument was communicated by \name{Per Alexandersson} in 2026.

\subsection[toeplitzPentadiagonal]{Balanced pentadiagonal symbols}

Suppose $b_{-2},\dotsc,b_2$ are real and $b_{-2}b_2>0$.
Replacing $b(z)$ by $(b(\rho z)-b_0)/c$, where
\begin{align*}
 \rho&=(b_{-2}/b_2)^{1/4}>0,\\
 c&=\operatorname{sgn}(b_2)\sqrt{b_{-2}b_2},
\end{align*}
preserves reality of the finite and limiting spectra: the substitution is
a diagonal similarity, followed by a real shift and nonzero real scaling.
We may therefore use the normalized symbol
\[
 b(z)=z^2+Az+Bz^{-1}+z^{-2},\qquad A,B\in\setR.
\]

\begin{proposition}
For this symbol the following conditions are equivalent:
\begin{enumerate}
\item $P_n$ is real-rooted for every $n\geq1$;
\item $P_n$ is real-rooted for every sufficiently large $n$;
\item $\Lambda(b)\subseteq\setR$;
\item $A=B$, or $(AB-16)^3\geq27(A^2-B^2)^2$.
\end{enumerate}
\end{proposition}

\begin{proof*}
The first two implications follow from the definition of $\Lambda(b)$.
If $A=B$, every section is real symmetric. Assume henceforth that $A\ne B$.
We first show that a nonreal critical point $w$ gives a nonreal point
$b(w)\in\Lambda(b)$. Put
\[
 q_t(z)=z^4+Az^3-tz^2+Bz+1,\qquad
 g(z)=z^3b'(z)=2z^4+Az^3-Bz-2.
\]
If $t=b(w)$ were real, the double root $w$ of $q_t$ would imply
$q_t(z)=(z-w)^2(z-\bar w)^2$. Its constant coefficient gives $|w|=1$,
and comparison of its cubic and linear coefficients gives $A=B$.
Thus $t$ is nonreal.

At this critical value,
\[
 q_t(z)=(z-w)^2\bigl(z^2+(A+2w)z+w^{-2}\bigr).
\]
The remaining roots divided by $w$ are the roots $y_1,y_2$ of
$y^2+py+d$, where $p=2+A/w$ and $d=w^{-4}$.
Write $R=|w|^2$ and $x=\operatorname{Re}w$.
Taking real and imaginary parts of $g(w)=0$, for $x\ne0$, gives
\[
 A=\frac{R-R^{-1}}x-4x,\qquad
 B=-\frac{R^2-1}x-\frac{4x}R.
\]
Also $p-\bar p d=(B-A/R)/w^3$. If $R\ne1$, these identities imply
\[
 \frac{|p-\bar p d|}{|1-|d|^2|}=\frac{\sqrt R}{|x|}>1.
\]
Now use the elementary identity
\[
 (1-|d|^2)^2-|p-\bar p d|^2
 =(1-|y_1|^2)(1-|y_2|^2)|1-y_1\bar y_2|^2.
\]
Its left side is negative, so exactly one $y_j$ lies inside the unit circle
and the other lies outside. The case $R=1$, $x\ne0$ would give $A=B$.
If $x=0$, criticality gives $w=\pm i$, $B=-A$, and the same left side
is $-4A^2<0$. Thus the double root occupies the middle two modulus ranks:
\[
 |z_1|<|z_2|=|z_3|<|z_4|.
\]
The Schmidt--Spitzer theorem gives $b(w)\in\Lambda(b)$.
Consequently, a real limiting spectrum forces all four roots of $g$ to be real.

Conversely, suppose all four critical points are real. Their product is $-1$,
so three have one sign and one has the other, counting multiplicity.
Replacing $z$ by $-z$ is an alternating diagonal similarity; we may assume
three critical points are negative. On the negative axis, $b(z)$ tends to
$+\infty$ at both ends. If these critical points are distinct, the two local
minima and intervening maximum give a level $t_0$ with four negative
solutions of $b(z)=t_0$. A double negative critical point $w$ instead gives
a triple root of $q_{b(w)}$, whose remaining root is $w^{-3}<0$;
a triple critical point gives a quadruple root. In all cases,
\[
 b(z)-t_0=z^{-2}\prod_{j=1}^4(z+\alpha_j),\qquad \alpha_j>0.
\]
Each lower bidiagonal matrix $T_{n+2}(z+\alpha_j)$ is
\hyperref[totallyNonnegative]{totally non-negative}.
The submatrix of their product in rows $3,\dotsc,n+2$ and columns
$1,\dotsc,n$ is $T_n(b)-t_0I$, hence is totally non-negative.
By the \hyperref[schurJacobiTrudi]{dual Jacobi--Trudi identity},
\[
 \det(T_n(b)-t_0I)=\schurS_{(n,n)}(\alpha_1,\dotsc,\alpha_4)>0.
\]
Its first upper and lower diagonals are positive. The oscillatory-matrix
criterion therefore gives positive simple eigenvalues, as in the valid
total-nonnegativity argument of \cite{ShapiroStampach2023Correction}.
Adding back $t_0I$ proves reality of every finite spectrum.

Finally,
\[
 \operatorname{Disc}(g)=4\bigl((AB-16)^3-27(A^2-B^2)^2\bigr).
\]
Since $g(0)=-2$ and $g$ is positive at both real infinities, it always
has a positive and a negative real root. For nonzero discriminant,
positive discriminant is therefore equivalent to four real roots.
At zero discriminant, a double real root and a nonreal conjugate pair,
or two conjugate double roots, would have positive product, contrary to $-1$.
This proves the parameter criterion, including its boundary.
\end{proof*}

\section[toeplitzReciprocalRecurrence]{Toeplitz determinants and reciprocal roots}

For the \hyperref[toeplitzPentadiagonal]{balanced pentadiagonal symbol}
\[
 b(z)=z^2+Az+Bz^{-1}+z^{-2},\qquad A,B\in\setR,
\]
put $P_n(t)=\det(tI-T_n(b))$ and $P_0=1$.
Let $z_1,\dotsc,z_4$ be the roots of
$q_t(z)=z^4+Az^3-tz^2+Bz+1$.
The \hyperref[schurJacobiTrudi]{dual Jacobi--Trudi identity} gives
$P_n(t)=(-1)^n\schurS_{(n,n)}(z_1,\dotsc,z_4)$.
The Toeplitz determinant formula and Schur recurrence in
\cite{Alexandersson2012Toeplitz} therefore give the six characteristic roots
\[
 \lambda_{ij}=-z_i z_j,\qquad 1\leq i<j\leq4.
\]
For distinct $z_i$, the corresponding expansion is
\[
 P_n(t)=\sum_{1\leq i<j\leq4}
 \frac{(z_i z_j)^2(-z_i z_j)^n}
 {\prod_{k\notin\{i,j\}}(z_i-z_k)(z_j-z_k)}.
\]
Thus the amplitudes are nonzero away from collisions of the symbol roots.
The characteristic polynomial is
\begin{align*}
 D(u)&=\prod_{i<j}(u+z_i z_j)\\
 &=u^6-tu^5+(AB-1)u^4+(2t+A^2+B^2)u^3\\
 &\qquad +(AB-1)u^2-tu+1.
\end{align*}
Equivalently, summing the expansion gives
\[
 G(u)\coloneqq\sum_{n\geq0}P_n(t)u^n=\frac{1-u^2}{D(u)}.
\]
These polynomial identities extend to repeated symbol roots by continuity.
In particular, $D$ specifies an order-six recurrence for $n\geq6$,
with initial polynomials determined by this generating function;
special parameters can reduce the minimal order.

Since $z_1z_2z_3z_4=1$, complementary pairs of rates multiply to $1$.
Writing $v=u+u^{-1}$ reduces $D(u)=0$ to
\[
 (v^2-4)(v-t)+ABv+A^2+B^2=0.
\]
Moreover, the rational generating function satisfies $G(1/u)=-u^4G(u)$.
Extending the order-six recurrence backwards gives
$P_{-1}=P_{-2}=P_{-3}=0$ and $P_{-4-n}=P_n$ for $n\geq0$.
Indeed, the expansion of $G$ at infinity is
$-\sum_{n<0}P_nu^n=-u^{-4}\sum_{n\geq0}P_nu^{-n}$.

Order the symbol roots by modulus. The two largest pair-product moduli
are $|z_3z_4|$ and $|z_2z_4|$, so a tie for dominance occurs precisely
when $|z_2|=|z_3|$. On analytic domains satisfying the
\hyperref[toeplitzBkwTheorem]{BKW nondegeneracy hypotheses}, the expansion
above explains this equal-modulus locus as a limiting zero set.
The \hyperref[toeplitzSpectra]{Schmidt--Spitzer theorem} supplies the
global statement, including collisions and exceptional parameter choices.
Reciprocal pairing gives the same locus for the backward recurrence.
For positive outer product, the
\hyperref[toeplitzPentadiagonal]{critical-point and total-nonnegativity proof}
shows that a real limiting spectrum forces every $P_n$ to be real-rooted.

\section[toeplitzHessenberg]{Five-term Hessenberg recurrences}

Consider
\[
 P_n(t)=tP_{n-1}(t)+\alpha P_{n-2}(t)+\beta P_{n-3}(t)+\gamma P_{n-4}(t),
 \qquad P_0=1,\quad P_{-1}=P_{-2}=P_{-3}=0,
\]
where $\alpha,\beta,\gamma\in\setR$ and $\gamma\ne0$.
These are characteristic polynomials of the Toeplitz sections with symbol
\[
 b(z)=-z^{-1}+\alpha z-\beta z^2+\gamma z^3,\qquad
 G(z)=z^2b'(z)=3\gamma z^4-2\beta z^3+\alpha z^2+1.
\]
The matrix is \defin{lower Hessenberg}: entries above the first
superdiagonal vanish. It has one upper and three lower diagonals, unlike the
\hyperref[toeplitzPentadiagonal]{balanced pentadiagonal case}.

\begin{theorem}[Sufficiency in \cite[Thm.~5]{Ndikubwayo2022Toeplitz}]
Every $P_n$ is real-rooted if either of the following holds:
\begin{enumerate}
\item[(a)] $G$ has four real roots, counted with multiplicity;
\item[(b)] $G$ has two distinct real roots of opposite signs and a nonreal
conjugate pair, and
\[
 24\alpha\gamma-5\beta^2\geq0,\qquad
 64\gamma^3+(4\alpha\gamma-\beta^2)^2\geq0.
\]
\end{enumerate}
\end{theorem}

The published proof invokes the
\hyperref[toeplitzSpectra]{Jordan-curve implication that fails in general}.
We give direct proofs of these sufficiency statements below.
The coefficient $24\alpha\gamma-5\beta^2$ is the published condition;
the earlier preprint has a different inequality.

\subsection[toeplitzRealCriticalHessenberg]{Four real critical points}

We first prove case (a). Replacing $z$ by $-z$ gives the equivalent symbol
\[
 f(z)=z^{-1}+az+dz^2-\gamma z^3,\qquad a=-\alpha,\quad d=-\beta.
\]

\begin{proof*}[The case $\gamma<0$]
Put $c=-\gamma>0$. The critical quartic $G(-z)$ has root product
$-1/(3c)$, so three critical points have one sign and one has the other.
Replacing $f(z)$ by $-f(-z)$, if needed, preserves reality of spectra and
makes three critical points negative.
On the negative axis $f$ tends to $-\infty$ at both ends.
If its three critical points there are distinct, two local maxima and the
intervening minimum give a level $t_0$ with four negative solutions.
At a double critical point $w<0$, the kernel has a triple root $w$;
its remaining root is $1/(cw^3)<0$. A triple critical point gives a
quadruple root. Thus in every case
\[
 f(z)-t_0=cz^{-1}\prod_{j=1}^4(z+\delta_j),\qquad \delta_j>0.
\]
The bidiagonal-product argument from the
\hyperref[toeplitzPentadiagonal]{balanced case}, now selecting rows
$2,\dotsc,n+1$ and columns $1,\dotsc,n$, makes $T_n(f)-t_0I$
totally non-negative. Its determinant is
$c^n\schurS_{(n,n,n)}(\delta_1,\dotsc,\delta_4)>0$, and its first
upper and lower diagonals are positive. The oscillatory-matrix criterion
gives real simple eigenvalues.
\end{proof*}

\begin{proof*}[The case $\gamma>0$]
Put $g=\gamma>0$, so that
\[
 f(z)=z^{-1}+az+dz^2-gz^3.
\]
Then
\[
 \sum_{n\geq0}P_n(t)u^n
 =\frac1{1-tu+au^2+du^3-gu^4}.
\]
The polynomial $-z^2f'(z)=1-az^2-2dz^3+3gz^4=G(-z)$
has four nonzero real roots. Their product is positive and the sum of
their reciprocals is zero, so two are positive and two negative.
Also $a>0$: the second elementary symmetric function of the reciprocal
roots equals $-\frac12$ times the sum of their squares, and equals $-a$.

Choose the smallest positive root $R$ of
\[
 H_*(r)=1-ar^2+2|d|r^3+3gr^4.
\]
Such a root exists because $H_*$ is one of $G(r),G(-r)$.
Since $H_*(0)>0$, its derivative at its first root is nonpositive.
Using $H_*(R)=0$ in $RH_*'(R)\leq0$ gives
\[
 |d|R^3+3gR^4\leq1. \tag{1}
\]

For $-1<x<1$, set
\[
 H_x(r)=1-ar^2-2dxr^3+g(4x^2-1)r^4.
\]
We have $H_x(R)<0$, since
\[
 H_x(R)-H_*(R)=-2(dx+|d|)R^3+4g(x^2-1)R^4<0.
\]
Moreover, on $0<r\leq R$,
\[
 \left(\frac{H_x(r)}{r^2}\right)''
 =\frac6{r^4}+2g(4x^2-1)>0
\]
by (1), which implies $R^4\leq1/(3g)$.
Since $H_x(r)/r^2$ tends to positive infinity at zero and is negative
at $R$, strict convexity gives exactly one zero $r(x)$ in $(0,R)$,
with negative derivative there. This zero varies continuously with $x$.

For $0<\theta<\pi$, take $x=\cos\theta$, $r=r(x)$, and
$t(\theta)=f(re^{i\theta})$. The equation $H_x(r)=0$ says exactly
that $t(\theta)$ is real. The two other roots of
$1-tz+az^2+dz^3-gz^4$ are the roots of
\[
 L(z)=z^2+(2rx-d/g)z-\frac1{gr^2}.
\]
They have opposite signs. In fact both lie outside $[-r,r]$, since
\[
 gr^2 L(\pm r)
 =gr^4(1\pm2x)\mp dr^3-1
 \leq3gr^4+|d|r^3-1<0.
\]
Write them as $s>r$ and $\ell<-r$, and put
$\eta=r/s\in(0,1)$, $\zeta=r/\ell\in(-1,0)$.

The map $\theta\mapsto t(\theta)$ is continuous and injective:
the quartic at a fixed real $t$ has just one nonreal conjugate pair,
which uniquely determines $\theta\in(0,\pi)$. Thus it is strictly monotone.

It remains to count finite-size zeros. Define
\[
 U_j(x)=\frac{\sin((j+1)\theta)}{\sin\theta},\qquad
 S_n(v)=\sum_{j=0}^n v^jU_{n-j}(x).
\]
The generating function factorization gives
\[
 r^nP_n(t(\theta))
 =\frac{\eta S_n(\eta)-\zeta S_n(\zeta)}{\eta-\zeta}.
\]
A finite geometric-sum identity gives
\[
 S_n(v)=\frac{U_n(x)-vU_{n+1}(x)+v^{n+2}}{1-2vx+v^2}.
\]
At $\theta_k=k\pi/(n+2)$, for $k=1,\dotsc,n+1$,
$U_{n+1}=0$ and $U_n=(-1)^{k+1}$. Since $|v|<1$,
both $S_n(\eta)$ and $S_n(\zeta)$ have this sign.
The weights $\eta$ and $-\zeta$ are positive, so
\[
 \operatorname{sgn}P_n(t(\theta_k))=(-1)^{k+1}.
\]
Strict monotonicity provides $n$ disjoint real intervals with sign changes.
Since $P_n$ is monic of degree $n$, all its roots are real and simple.
The argument includes repeated roots of $G$; only the strict interior
inequalities $|x|<1$ and $r<R$ were used. 

\end{proof*}

\subsection[toeplitzQuarticKernel]{A quartic-kernel sufficient condition}

For case (b), we prove a sufficient condition without a critical-root
hypothesis. The two conjugate root pairs have different moduli, and a
residue estimate gives alternating signs at suitable real values of $t$.

\begin{proposition}
Suppose $a,d\in\setR$ and $c>0$, and define the monic polynomials $Q_n$ by
\[
 \sum_{n\geq0}Q_n(t)u^n=\frac1{1-tu+au^2+du^3+cu^4}.
\]
If $a\geq2\sqrt c+d^2/(4c)$, then every $Q_n$ is real-rooted.
If the inequality is strict, all its roots are simple.

\end{proposition}

\begin{proof*}
Replacing $u$ by $c^{-1/4}u$ reduces to $c=1$,
with parameters $a/\sqrt c,d/c^{3/4}$ and variable $t/c^{1/4}$.
We first assume $a>2+d^2/4$. Put
\[
 f(z)=z^{-1}+az+dz^2+z^3.
\]
For $-1\leq x\leq1$, consider
\[
 H_x(r)=1-ar^2-2dxr^3+(1-4x^2)r^4.
\]
We have
\begin{align*}
 H_x(1)&=2-a+d^2/4-(2x+d/2)^2<0,\\
 (H_x(r)/r^2)''&=6/r^4+2(1-4x^2)>0\quad(0<r<1).
\end{align*}
Since $H_x(r)/r^2$ tends to positive infinity at zero, strict convexity
gives a unique zero $r(x)\in(0,1)$, with negative derivative there.
It depends continuously on $x$, including the endpoints.

Take $x=\cos\theta$, $0<\theta<\pi$, and $r=r(x)$.
Then $t(\theta)=f(re^{i\theta})$ is real, since
$\operatorname{Im}f(re^{i\theta})=-H_x(r)\sin\theta/r$.
The kernel $1-tz+az^2+dz^3+z^4$ has roots $re^{\pm i\theta}$.
Write $\eta,\zeta$ for $r$ divided by its other two roots, and set
$\rho=r^2$. Comparing the constant and cubic coefficients gives
\[
 \eta\zeta=\rho^2,\qquad
 \eta+\zeta=-dr^3-2\rho^2x.
\]
The identity $H_x(r)=0$ and $a>2+d^2/4$ imply
\[
 (1-\rho)^2>
 \left(\frac{dr^3}{2\rho}+2\rho x\right)^2.
\]
Define $y=-dr^3/(2\rho)-\rho x$. Thus
\[
 |y-\rho x|<1-\rho,\qquad |y|<1.
\]
Consequently $\eta,\zeta=\rho e^{\pm i\phi}$, where
$\phi=\arccos y\in(0,\pi)$. Moreover,
\[
 |y|\leq\rho|x|+1-\rho
 \quad\Longrightarrow\quad
 \sin\phi\geq\rho\sin\theta.
\]
Indeed, subtracting $\rho^2(1-x^2)$ from
$1-(\rho|x|+1-\rho)^2$ gives
$2\rho(1-\rho)(1-|x|)\geq0$.

The two smaller kernel roots have modulus $r<1$, and the other two
have modulus $r/\rho=1/r>1$. Thus $t(\theta)$ is injective:
its value determines its unique smaller conjugate pair and hence $\theta$.
It is therefore strictly monotone.

The coefficient of $v^n$ in
\[
 \frac1{(1-e^{i\theta}v)(1-e^{-i\theta}v)
 (1-\rho e^{i\phi}v)(1-\rho e^{-i\phi}v)}
\]
is $r^nQ_n(t(\theta))$. For four distinct numbers $\lambda_j$,
this coefficient equals
$\sum_j\lambda_j^{n+3}/\prod_{k\ne j}(\lambda_j-\lambda_k)$.
All four numbers here are distinct. Put
\begin{align*}
 A&=(e^{i\theta}-\rho e^{i\phi})(e^{i\theta}-\rho e^{-i\phi}),\\
 \psi&=\arg\bigl((1+\rho^2)x-2\rho y
                 +i(1-\rho^2)\sin\theta\bigr).
\end{align*}
The argument lies in $(0,\pi)$, and $A=e^{i\theta}|A|e^{i\psi}$.
Combining the residues of the first conjugate pair gives
\begin{align*}
 r^nQ_n(t(\theta))
 &=\frac{\sin((n+2)\theta-\psi)}{\sin\theta\,|A|}+E_n(\theta),\\
 |E_n(\theta)|&\leq\frac{\rho^{n+2}}{\sin\phi\,|A|}\\
 &\leq\frac{\rho^{n+1}}{\sin\theta\,|A|}.
\end{align*}
For the error bound, the absolute product of the two cross-pair differences
at $\rho e^{i\phi}$ equals $|A|$; its within-pair difference has
absolute value $2\rho\sin\phi$. The last inequality uses $\sin\phi\geq\rho\sin\theta$.

The function $\psi$ extends continuously to $[0,\pi]$, with
$\psi(0)=0$ and $\psi(\pi)=\pi$. At the first endpoint its real
part is $1+\rho^2-2\rho y>0$, and at the second it is
$-1-\rho^2-2\rho y<0$. Hence the continuous phase
$(n+2)\theta-\psi$ runs from $0$ to $(n+1)\pi$.
We can select successively ordered angles where it equals
$(k+1/2)\pi$, for $k=0,\dotsc,n$. At these angles the residue estimate, with
$\rho<1$, shows that $Q_n(t(\theta))$ has sign $(-1)^k$.
Strict monotonicity of $t(\theta)$ now gives $n$ distinct real roots.
Finally, replacing $a$ by $a+\varepsilon$ and taking
$\varepsilon\downarrow0$ proves the weak-inequality case, since
the fixed-degree monic real-rooted polynomials form a closed set.
\end{proof*}

\subsection[toeplitzPublishedHessenbergRegion]{The published two-real-critical-point region}

In case (b) of \cite[Thm.~5]{Ndikubwayo2022Toeplitz}, $G$ has exactly two distinct real roots of opposite
signs and one nonreal conjugate pair. Its root product is negative, so
$\gamma<0$. Set
\[
 c=-\gamma>0,\qquad a=-\alpha,\qquad d=-\beta.
\]
The published inequalities are
\[
 24ac\geq5d^2,\qquad (4ac-d^2)^2\geq64c^3.
\]
After the preceding normalization, these read
$24a\geq5d^2$ and $|4a-d^2|\geq8$.
We show that the negative branch $4a-d^2\leq-8$ is impossible.

Put $v=d^2$. On that branch,
\[
 4a+8\leq v\leq24a/5,\qquad a\geq10.
\]
The critical quartic is $1-az^2-2dz^3-3z^4$, up to a real
nonzero scaling of its argument. Its discriminant is $16D(a,v)$, where
\[
 D(a,v)=-27v^2+(a^3+108a)v-3a^4-72a^2-432.
\]
For fixed $a$, this is concave in $v$, and its endpoint values are
\begin{align*}
 D(a,4a+8)&=(a-10)(a+6)^3,\\
 D(a,24a/5)&=\frac9{25}(a^2-100)(5a^2+12).
\end{align*}
Both are nonnegative for $a\geq10$. This contradicts the negative
discriminant of a real quartic with two simple real roots and a nonreal
conjugate pair. Therefore $4a-d^2\geq8$, which is precisely
$a\geq2+d^2/4$. The proposition applies.

In the original parameters, the sufficient condition used is
\[
 -\alpha\geq2\sqrt{-\gamma}+\frac{\beta^2}{-4\gamma}.
\]
This proves case (b) of the sufficiency theorem.
